\section*{Evaluasi}
\addcontentsline{toc}{section}{Evaluasi}

\begin{enumerate}[label=\arabic*.,leftmargin=*]
  \item Uraikan tiga alasan mengapa pemahaman teknik kompilasi berguna
        bagi seorang insinyur perangkat lunak.
  \item Buat tabel yang memetakan empat CPMK ke minimal dua bab inti
        masing-masing.
  \item Jelaskan perbedaan peran Bab I, Bab II, dan Bab XVII dalam
        struktur buku.
  \item Diberikan potongan program Subset-C, tuliskan urutan token dan
        jelaskan aturan yang Anda pakai.
  \item Rancang rencana belajar empat belas minggu yang memuat bacaan,
        praktikum, dan evaluasi mandiri.
\end{enumerate}

\subsection*{Kunci dan Rubrik Singkat}
\begin{itemize}[leftmargin=*]
  \item Soal 1--3: ketepatan konsep (60\%), kelengkapan (25\%),
        kejelasan bahasa (15\%).
  \item Soal 4: kebenaran token (70\%), penjelasan aturan (30\%).
  \item Soal 5: kelayakan jadwal (50\%), keterukuran target (30\%),
        keseimbangan teori-praktik (20\%).
\end{itemize}
